\documentclass[10pt,a4paper]{amsart}
\usepackage{amsmath,amssymb,amsthm,graphicx}
% For commutative diagrams we use Xy-pic
\usepackage[all]{xy}

% We use 2cell for 2-commutative diagrams.
\xyoption{2cell}
% Theorem environments.
%
\theoremstyle{plain}
\newtheorem{theorem}[subsection]{Theorem}
\newtheorem{proposition}[subsection]{Proposition}
\newtheorem{lemma}[subsection]{Lemma}

\theoremstyle{definition}
\newtheorem{definition}[subsection]{Definition}
\newtheorem{example}[subsection]{Example}
\newtheorem{exercise}[subsection]{Exercise}
\newtheorem{situation}[subsection]{Situation}

\theoremstyle{remark}
\newtheorem{remark}[subsection]{Remark}
\newtheorem{remarks}[subsection]{Remarks}

\numberwithin{equation}{subsection}

% Macros
%
\def\lim{\mathop{\mathrm{lim}}\nolimits}
\def\colim{\mathop{\mathrm{colim}}\nolimits}
\def\Spec{\mathop{\mathrm{Spec}}}
\def\Hom{\mathop{\mathrm{Hom}}\nolimits}
\def\Ext{\mathop{\mathrm{Ext}}\nolimits}
\def\SheafHom{\mathop{\mathcal{H}\!\mathit{om}}\nolimits}
\def\SheafExt{\mathop{\mathcal{E}\!\mathit{xt}}\nolimits}
\def\Sch{\mathit{Sch}}
\def\Mor{\mathop{\mathrm{Mor}}\nolimits}
\def\Ob{\mathop{\mathrm{Ob}}\nolimits}
\def\Sh{\mathop{\mathit{Sh}}\nolimits}
\def\NL{\mathop{N\!L}\nolimits}
\def\CH{\mathop{\mathrm{CH}}\nolimits}
\def\proetale{{pro\text{-}\acute{e}tale}}
\def\etale{{\acute{e}tale}}
\def\QCoh{\mathit{QCoh}}
\def\Ker{\mathop{\mathrm{Ker}}}
\def\Im{\mathop{\mathrm{Im}}}
\def\Coker{\mathop{\mathrm{Coker}}}
\def\Coim{\mathop{\mathrm{Coim}}}

% Boxtimes
%
\DeclareMathSymbol{\boxtimes}{\mathbin}{AMSa}{"02}

%
% Macros for moduli stacks/spaces
%
\def\QCohstack{\mathcal{QC}\!\mathit{oh}}
\def\Cohstack{\mathcal{C}\!\mathit{oh}}
\def\Spacesstack{\mathcal{S}\!\mathit{paces}}
\def\Quotfunctor{\mathrm{Quot}}
\def\Hilbfunctor{\mathrm{Hilb}}
\def\Curvesstack{\mathcal{C}\!\mathit{urves}}
\def\Polarizedstack{\mathcal{P}\!\mathit{olarized}}
\def\Complexesstack{\mathcal{C}\!\mathit{omplexes}}
% \Pic is the operator that assigns to X its picard group, usage \Pic(X)
% \Picardstack_{X/B} denotes the Picard stack of X over B
% \Picardfunctor_{X/B} denotes the Picard functor of X over B
\def\Pic{\mathop{\mathrm{Pic}}\nolimits}
\def\Picardstack{\mathcal{P}\!\mathit{ic}}
\def\Picardfunctor{\mathrm{Pic}}
\def\Deformationcategory{\mathcal{D}\!\mathit{ef}}
\usepackage{cleveref}
\crefname{originallemma}{original lemma}{original lemmas}
\Crefname{originallemma}{Original lemma}{Original lemmas}
\makeatletter
\newlabel{algebra-lemma-quotient-strict-henselization}{{05WS}{source}{}{sourcepointer-1}{}}
\newlabel{algebra-lemma-quotient-strict-henselization@cref}{{[originallemma][0][]05WS}{[1][0][]source}{}{}{}}
\makeatother
\crefname{originalreference}{original reference}{original references}
\Crefname{originalreference}{Original reference}{Original references}
\makeatletter
\newlabel{lemma-gabber-criterion}{{0GY0}{source}{}{sourcepointer-2}{}}
\newlabel{lemma-gabber-criterion@cref}{{[originalreference][0][]0GY0}{[1][0][]source}{}{}{}}
\makeatother
\crefname{originalreference}{original reference}{original references}
\Crefname{originalreference}{Original reference}{Original references}
\makeatletter
\newlabel{lemma-gabber-for-h1-absolutely-algebraically-closed}{{09ZB}{source}{}{sourcepointer-3}{}}
\newlabel{lemma-gabber-for-h1-absolutely-algebraically-closed@cref}{{[originalreference][0][]09ZB}{[1][0][]source}{}{}{}}
\makeatother
\crefname{originalreference}{original reference}{original references}
\Crefname{originalreference}{Original reference}{Original references}
\makeatletter
\newlabel{lemma-normal-scheme-with-alg-closed-function-field}{{09Z9}{source}{}{sourcepointer-4}{}}
\newlabel{lemma-normal-scheme-with-alg-closed-function-field@cref}{{[originalreference][0][]09Z9}{[1][0][]source}{}{}{}}
\makeatother
\title{Gabber vanishing for a closed subscheme of a normal integral affine scheme with separably closed function field}
\author{The Stacks Project Authors}
\date{stacks-project:etale-cohomology.tex:local-snapshot-20261011; GNU FDL 1.2 or later}
\begin{document}
\maketitle
\noindent\textit{Editorial note.} All delivered mathematical text is the author's own, copied from the pinned Stacks Project chapter file etale-cohomology.tex; no statement or proof step was rewritten, reordered or removed.\par\medskip
\noindent\textit{Editorial note.} The counted proof is the author's complete argument: the proof is written out in full: the presentation of the closed subscheme, the strict henselianity of its local rings, the h\textasciicircum{}1 vanishing on principal opens with the lift argument, and the application of Gabber's criterion.\par\medskip
\noindent\textit{Editorial note.} The label lemma-gabber-for-absolutely-algebraically-closed names the counted statement, and the proof environment immediately following it in the source is delivered verbatim as its proof.\par\medskip
\noindent\textit{Editorial note.} The counted statement is carried by the author's own lemma environment in the chapter source.\par\medskip
\noindent\textit{Editorial note.} The statement and proof are copied from the pinned chapter file; the four cited inputs are kept as external pointers with their canonical Stacks tags.\par\medskip
\noindent\textit{Allowed standard prerequisites (editorial).}\begin{itemize}
\item Strict henselianity of the local rings of the quotient. Same chapter, the structure lemma for affine normal schemes with separably closed function field, together with Algebra, the lemma that quotients of strict henselizations are strictly henselian. Supplies the strictly henselian local rings of the closed subscheme.
\item The h\textasciicircum{}1 vanishing for absolutely algebraically closed normal domains. Same chapter, the lemma stating the vanishing of H\textasciicircum{}1 with finite constant coefficients over a normal domain with separably closed fraction field. Supplies the vanishing on the principal opens of the closed subscheme.
\item Gabber's criterion for the vanishing of cohomology. Same chapter, the criterion reducing the vanishing of all higher cohomology on a scheme to conditions on its local rings and principal opens. Supplies the final step of the proof.
\end{itemize}\medskip
\ifdefined\hypertarget\hypertarget{sourcepointer-1}{}\fi
\noindent\textit{Source reference.} Original lemma 05WS: sources/stacks-chow/tags.txt, source lines 7596--7596. Prerequisite pointer invoked at proof line(s) 14456; it names a standard fact of a sibling chapter that the counted argument uses, and every step of the counted proof built on it is retained verbatim in the delivered body.\par\medskip
\ifdefined\hypertarget\hypertarget{sourcepointer-2}{}\fi
\noindent\textit{Source reference.} Original reference 0GY0: sources/stacks-etale-cohomology/pinned-original-source.tex, source lines 14228--14228. Prerequisite pointer invoked at proof line(s) 14462; it names a standard fact of the same chapter that the counted argument uses, and every step of the counted proof built on it is retained verbatim in the delivered body.\par\medskip
\ifdefined\hypertarget\hypertarget{sourcepointer-3}{}\fi
\noindent\textit{Source reference.} Original reference 09ZB: sources/stacks-etale-cohomology/pinned-original-source.tex, source lines 14425--14425. Prerequisite pointer invoked at proof line(s) 14461; it names a standard fact of the same chapter that the counted argument uses, and every step of the counted proof built on it is retained verbatim in the delivered body.\par\medskip
\ifdefined\hypertarget\hypertarget{sourcepointer-4}{}\fi
\noindent\textit{Source reference.} Original reference 09Z9: sources/stacks-etale-cohomology/pinned-original-source.tex, source lines 14312--14312. Prerequisite pointer invoked at proof line(s) 14455; it names a standard fact of the same chapter that the counted argument uses, and every step of the counted proof built on it is retained verbatim in the delivered body.\par\medskip
\begin{lemma}
\label{lemma-gabber-for-absolutely-algebraically-closed}
Let $X$ be a normal integral affine scheme with separably closed
function field. Let $Z \subset X$ be a closed subscheme.
For any finite abelian group $M$ we have
$H^q_\etale(Z, \underline{M}) = 0$ for $q \geq 1$.
\end{lemma}

\begin{proof}
Write $X = \Spec(R)$ and $Z = \Spec(R')$ so that we have a surjection
of rings $R \to R'$. All local rings of $R'$ are strictly henselian by
Lemma \ref{lemma-normal-scheme-with-alg-closed-function-field} and
Algebra, Lemma \ref{algebra-lemma-quotient-strict-henselization}.
Furthermore, we see that for any $f' \in R'$ there is a surjection
$R_f \to R'_{f'}$ where $f \in R$ is a lift of $f'$.
Since $R_f$ is a normal domain with separably closed fraction field
we see that $H^1_\etale(D(f'), \underline{M}) = 0$ by
Lemma \ref{lemma-gabber-for-h1-absolutely-algebraically-closed}.
Thus we may apply Lemma \ref{lemma-gabber-criterion} to $Z = \Spec(R')$
to conclude.
\end{proof}

\end{document}
